% Учебный пример для проверки загрузки; это не научная публикация.
\section{Экспоненциальное забывание: учебный пример}
Рассмотрим состояние самостоятельного успеха x(t) при отсутствии практики.
Уравнение и начальное условие:
$$\dot x=-\lambda x,\qquad x(0)=X.$$
Здесь X – начальная вероятность, lambda – положительная скорость забывания,
t – время в условных учебных периодах. Решение:
$$x(t)=X e^{-\lambda t}.$$
Проверка: производная экспоненты равна минус lambda, умноженному на ту же экспоненту.
При X=0,65, lambda=0,22 и t=1 получаем приблизительно 0,52164.
Это синтетическая иллюстрация, а не оценка забывания реального студента.

\section{Вопросы к лектору}
Почему решение положительно? Почему lambda не равна доле потерянных знаний
за весь курс? Как из двух точных самостоятельных вероятностей восстановить lambda?
Почему два бинарных ответа не позволяют подставить их в формулу отношения вероятностей?
